\documentclass[11pt]{article}
\usepackage{amsmath,amssymb,amsthm}
\usepackage[margin=1.1in]{geometry}
\usepackage{hyperref}
\newtheorem{theorem}{Theorem}[section]
\newtheorem{lemma}[theorem]{Lemma}
\newtheorem{proposition}[theorem]{Proposition}
\newtheorem{corollary}[theorem]{Corollary}
\newtheorem{definition}[theorem]{Definition}
\newtheorem{assumption}{Assumption}
\newtheorem{remark}[theorem]{Remark}
\newcommand{\Tr}{\operatorname{Tr}}
\newcommand{\G}{{\mathrm{SU}}(2)}
\DeclareMathOperator{\Ad}{Ad}
\title{\bf Exact second-order strong-coupling expansion of the spectral gap\\
of $\G$ Kogut--Susskind lattice gauge theory on the three-torus}
\author{Revised draft for external human review --- campaign YM-04.6\\
\small (verification artifacts accompany this manuscript; original YM-04 archive preserved unchanged)}
\date{\today}
\begin{document}
\maketitle

\begin{abstract}
We compute the first two nontrivial orders of the strong-coupling expansion of
the spectral gap of the pure $\G$ Kogut--Susskind Hamiltonian on the periodic
three-torus $\mathbb T_L^3$ of side $L$, in the dimensionless normalization
$H(\mu)=H_E+\mu H_M$ with $H_E=\sum_{\ell}C_2(\ell)$ and
$H_M=\sum_p(2-\Re\Tr U_p)$; the dictionary to the Kogut--Susskind coupling at
unit lattice spacing is $\mu=2g^{-6}$, so that in physical units
$\Delta_{\rm phys}(g,L)=\tfrac32 g^{2}-\tfrac{142}{105}\,g^{-10}+O(g^{-16})$. For every $L\ge 5$ the gap above the vacuum obeys
\[
\Delta(\mu,L)\;=\;3-\frac{71}{105}\,\mu^2+O_L(\mu^3),\qquad \mu\to 0^+,
\]
with an $L$-independent coefficient; the band head is the translation-invariant
one-plaquette excitation. At $L=4$ an additional winding sector enters the
$E=3$ eigenspace and the gap is governed by that sector with coefficient at
most $-16/21$. All group-theoretic coefficients are exact rational numbers
derived from Peter--Weyl orthogonality, $6j$ recoupling, and the
$\G$ Weingarten calculus; every step is machine-verifiable in exact rational
arithmetic. The statement is a fixed-volume perturbative theorem; no
uniformity in $L$ and no continuum limit are claimed. Claim hierarchy
(2026-11-09): the identification of the second-order coefficient with the
full spectral-gap expansion uses the channel-exhaustion selection rule R3'
of Appendix~G, which the main text states as an explicit \emph{Assumption}
pending human verification; all coefficient identities, the band-head
curvature and uniqueness analysis, and the conversion dictionary are
unconditional and machine-checked.
\end{abstract}

\section{Introduction and result}
Strong-coupling expansions of Hamiltonian lattice gauge theory go back to
Kogut and Susskind \cite{KS75} and to the glueball calculations of Kogut,
Sinclair and Susskind \cite{KSS76}; textbook accounts appear in
\cite{MM94,DZ83}. Within that framework we obtain an exact coefficient of the
second-order expansion of the spectral gap, together with a complete
finite-volume band analysis, in a normalization fixed once and for all from
the physical Hilbert space. The exact rational value and the finite-volume
band-head statement are the objects of interest; their status relative to the
published literature is discussed in Section~9 and in the accompanying
novelty ledger.

Let $\Lambda_L=(\mathbb Z/L\mathbb Z)^3$, $L\ge 4$, with link set $E$ and
plaquette set $\mathcal P$ ($|\mathcal P|=3L^3$). The physical Hilbert space
is $\mathcal H_0=\{f\in L^2(\G^{E},dU)\}$, the gauge-invariant subspace,
$dU$ the product Haar measure. The Hamiltonian is
\begin{equation}\label{eq:H}
H(\mu)\;=\;H_E+\mu H_M,\qquad
H_E=\sum_{\ell\in E}C_2(\ell),\qquad
H_M=\sum_{p\in\mathcal P}\bigl(2-\Re\Tr U_p\bigr),
\end{equation}
where $C_2(\ell)$ is the (left) Casimir on link $\ell$, acting as $j(j+1)$ on
the Peter--Weyl component of spin $j$, and $U_p$ is the ordered holonomy
around plaquette $p$. $H_M\ge 0$, $H_E\ge 0$; $H(\mu)$ is self-adjoint,
bounded below, with discrete spectrum at fixed $L$ ($H_E$ has compact
resolvent, $H_M$ is bounded). The vacuum $\Omega\equiv 1$ has $H(0)\Omega=0$;
for $\mu>0$ small it is the unique ground state (analytic perturbation
theory, Section~3).

\begin{assumption}[R3', channel exhaustion]\label{ass:r3prime}
Every intermediate state $n$ with $\langle n|H_M|p\rangle\neq 0$ that
contributes to the second-order term is of the recoupling type analyzed in
Appendices~A--B: writing $j_\ell(n)$ for the link spins of $n$ and $p,q$ for
the two plaquettes involved, $j_\ell(n)=\tfrac12$ on
$\mathrm{supp}(q)\,\Delta\,\mathrm{supp}(p)$, $j_\ell(n)\in\{0,1\}$ on
$\mathrm{supp}(p)\cap\mathrm{supp}(q)$, $j_\ell(n)=0$ outside
$\mathrm{supp}(p)\cup\mathrm{supp}(q)$; and the only nonvanishing overlap
contributions are the singlet and triplet channels of Appendix~B, with
weights $c_0^2=\tfrac14$, $c_1^2=\tfrac34$. The exact exhaustive form is
Theorem~\ref{thm:r3prime}; Appendix~G contains the claimed proof, offered
for human audit.
\end{assumption}

\begin{theorem}[main]\label{thm:main}
Under Assumption~\ref{ass:r3prime}, for every fixed $L\ge 5$ there exists
$\mu_0(L)>0$ such that for $0<\mu<\mu_0(L)$ the spectral gap of $H(\mu)$
above its ground state is
\[
\Delta(\mu,L)\;=\;3-\frac{71}{105}\,\mu^2+O_L(\mu^3),\qquad \mu\to0^+ .
\]
The coefficient $-71/105$ is independent of $L$. The second-order band
minimum is attained at crystal momentum $k=0$, in the uniform
orientation-symmetric one-plaquette mode, uniquely among band minima at
second order.
\end{theorem}

Without Assumption~\ref{ass:r3prime} the unconditional content is: the exact
rational value $-71/105$ of the second-order coefficient of the truncated
second-order expression over the Appendices~A--B channels (machine-checked),
the band-head curvature and Perron--Frobenius uniqueness, and every
conversion identity. Assumption~R3' enters exactly once: to identify that
channel sum with the second-order coefficient of the full resolvent
expansion.
\begin{remark}
On 2026-11-09 an external review declined to accept the Appendix~G
exhaustion argument as human-review grade. Rather than conceal the point,
this paper states R3' as a named assumption and records the full
adversarial-review trail in \texttt{evidence/REFEREE\_R3PRIME*.md}.
\end{remark}

\begin{corollary}[physical units]\label{cor:phys}
With $H_{\rm KS}=\frac{g^2}{2}\sum_\ell E^2(\ell)-\frac{1}{2g^4}\sum_p
(\Tr U_p+\Tr U_p^\dagger)$ at $a=1$ and $L\ge5$,
\[
\Delta_{\rm phys}(g,L)
=\frac{g^2}{2}\,\Delta\bigl(2g^{-6},L\bigr)
=\frac{3}{2}\,g^{2}-\frac{142}{105}\,g^{-10}+O_L(g^{-16})
=\frac{3}{2}\,g^{2}\Bigl(1-\frac{284}{315}\,g^{-12}+O_L(g^{-18})\Bigr),
\qquad g\to\infty .
\]
Every conversion factor is exact rational; a source quoting
$\mu=g^{-4}$ for the same operator differs by the factor $c=2$ and must be
rescaled by $c_2\mapsto c_2/c^{2}$ before comparison.
\end{corollary}

\begin{theorem}[$L=4$ exception]\label{thm:L4}
At $L=4$ the $E=3$ eigenspace of $H_E$ contains, besides the $192$
plaquette states, $48$ straight winding-loop states. The spectral gap is then
governed by the winding sector, whose second-order coefficient satisfies
$c_2(L{=}4)\le -16/21<-71/105$; the exact value of $c_2(L{=}4)$ is not
determined in this paper (Section~8).
\end{theorem}

Both theorems are fixed-volume statements. They say nothing about
$\inf_L\Delta(\mu,L)$ uniformity, and nothing about the continuum limit.
Section~9 states the precise scope boundary.

\section{Setup, conventions, and spectral preliminaries}\label{sec:setup}
\subsection{Hilbert space and gauge invariance}
$\mathcal H=L^2(\G^E,dU)$ with $dU=\prod_\ell dU_\ell$ the product of Haar
measures normalized to $1$. Gauge transformations act by
$(g\cdot f)(U)=f(\{\gamma_\ell^{-1}U_\ell\gamma_{\ell'}\})$ for vertex
functions $\gamma_x\in\G$; $\mathcal H_0$ is the invariant subspace. In the
Peter--Weyl decomposition
\[
L^2(\G)\;\cong\;\widehat{\bigoplus}_{j\in\frac12\mathbb N}
V_j\otimes V_j^*,\qquad C_2\big|_{V_j}=j(j+1),
\]
gauge invariance means: at each vertex $x$, the tensor product of the
$V_j$'s carried by links incident to $x$ contains the singlet, and the state
is the corresponding intertwiner contraction. A basis of $\mathcal H_0$ is
labelled by flux assignments $\{j_\ell\}$ plus intertwiners at vertices; the
$H_E$-eigenvalue is $\sum_\ell j_\ell(j_\ell+1)$.

\subsection{Normalization}
The primary object of this paper is the dimensionless family \eqref{eq:H}
with $\mu>0$ as the expansion parameter; $C_2(\ell)=E^2(\ell)$ on class
functions. Dictionary to the Kogut--Susskind Hamiltonian
$H_{\rm KS}=\frac{g^2}{2}\sum_\ell E^2(\ell)
-\frac{1}{2g^4}\sum_p(\Tr U_p+\Tr U_p^\dagger)$ at $a=1$: in units of
$g^2/2$,
\[
\frac{2}{g^2}H_{\rm KS}+\frac{4|\mathcal P|}{g^6}
=\sum_\ell C_2(\ell)+\frac{2}{g^6}\sum_p\bigl(2-\Re\Tr U_p\bigr),
\qquad\text{i.e.}\quad \mu=\frac{2}{g^6} .
\]
Some references (and the earlier notes of this campaign) absorb the factor of
$2$ into the coupling and quote $\mu=g^{-4}$ under the equally standard
convention $H_B=-g^{-4}\sum_p\Re\Tr U_p$. All results below are statements
about $H(\mu)$ itself; a change of convention $\mu\mapsto c\,\mu$ rescales
the coefficient as $c_2\mapsto c_2/c^2$. Conversion factors to
$\lambda=\tfrac12 g^{-4}$-type variables are tabulated in Appendix~D.

\subsection{Analyticity}
$H_M$ is bounded ($\|\Re\Tr U_p\|_\infty\le 2$), so $H(\mu)$ is a bounded
analytic family for $\mu\in\mathbb R$ in the norm-resolvent sense at fixed
$L$. The spectrum is discrete; simple eigenvalues and isolated clusters admit
convergent power series in $\mu$ (Kato). The vacuum cluster and the
$E=3$ cluster of $H(0)$ are separated by $3$, so for $\mu$ small the gap is
the bottom of the deformed $E=3$ cluster minus the deformed ground energy,
each analytic in $\mu$ on the real axis (the $E=3$ cluster is degenerate; its
splitting is handled by degenerate perturbation theory, which yields the
expansions of Theorem~\ref{thm:main} with $O_L(\mu^3)$ remainder).

\section{The $E=3$ eigenspace and Wilson-loop orthonormality}\label{sec:E3}
\begin{lemma}[flux classification]\label{lem:flux}
The $E=3$ eigenspace of $H_E$ on $\mathcal H_0$ is spanned by gauge-invariant
states carrying flux $j=\tfrac12$ on exactly four links and $0$ elsewhere.
For $L\ge 5$ the supporting links form a plaquette; for $L=4$ they form a
plaquette or a straight winding loop of length $4$.
\end{lemma}
\begin{proof}
$3=\sum_\ell j_\ell(j_\ell+1)$ with each $j_\ell(j_\ell+1)\in\{0,\tfrac34,2,
\tfrac{15}{4},\dots\}$ forces four links with $j=\tfrac12$ (since
$2+\tfrac34<3$ and $\tfrac{15}4>3$, and $3/4$ does not divide the remainder
$1$ otherwise). Gauge invariance requires the total flux at each vertex to
contain a singlet, so each vertex carries an even number ($0$ or $2$, as the
$j$'s are $\tfrac12$) of flux links: the support is a closed loop of length
$4$. On $\mathbb T_L^3$ a closed nearest-neighbor loop of length $4$ is a
plaquette, except when $L=4$ and the loop wraps once in a coordinate
direction (a straight winding loop).
\end{proof}

\begin{lemma}[singlet uniqueness and the Wilson-loop basis]\label{lem:singlet}
For a fixed closed loop $C=(\ell_1,\dots,\ell_4)$ with $j=\tfrac12$ on each
link, the space of gauge-invariant tensors is one-dimensional, spanned by
the Wilson loop function $\chi_{1/2}(U_C)=\Tr(U_{\ell_1}\cdots U_{\ell_4})$.
\end{lemma}
\begin{proof}
Each interior vertex of the loop pairs exactly two $V_{1/2}$ factors, whose
singlet space is one-dimensional ($V_{1/2}\otimes V_{1/2}\supset V_0$ with
multiplicity one, via the invariant $\varepsilon\in V_{1/2}\otimes V_{1/2}$).
The product of the four vertex singlets, propagated around the loop, is the
cyclic index contraction $\Tr(U_{\ell_1}\cdots U_{\ell_4})$; uniqueness of
the singlet at each vertex forces any intertwiner to be a scalar multiple of
it.
\end{proof}

\begin{lemma}[orthonormality]\label{lem:gram}
For loops $C,C'$ of length $4$ as in Lemma \ref{lem:flux},
\[
\bigl\langle \chi(U_C),\chi(U_{C'})\bigr\rangle_{L^2(\G^E)}
=\delta_{CC'} .
\]
\end{lemma}
\begin{proof}
$\overline{\chi(U_{C'})}=\chi(U_{C'}^\dagger)$, so the inner product is the
integral of a product carrying one $U$-type factor from $\chi(U_C)$ and, on
$C'$, one $\bar U$-type factor from $\overline{\chi(U_{C'})}$. If $C\ne C'$,
some link $\ell\in C\setminus C'$ carries a single $U$-type factor and no
$\bar U$-type factor; the remaining integrations produce a free matrix index
on $\ell$, and $\int U^a{}_b(U_\ell)\,dU_\ell=0$ because the spin-$\tfrac12$
matrix coefficient has no invariant vector. Hence the integral vanishes.
If $C=C'$, Peter--Weyl orthogonality
$\int U^a{}_b\,\overline{U^c{}_d}\,dU=\frac12\delta^a_d\delta^c_b$ is
applied link by link: each of the four links contributes a factor $\tfrac12$
and each of the four successive index sums a factor $2$, so the closed loop
evaluates to $2^4\cdot 2^{-4}=1$. The full index chain is spelled out in
Appendix~A and machine-verified by two independent exact methods (Weyl
integration and Weingarten) in \texttt{verify\_anchors.py} \cite{YM04code}.
\end{proof}
The proof of the diagonal case is spelled out with all index positions in
Appendix~A; it is also machine-verified by two independent exact methods
(Weyl integration and Weingarten) in \texttt{verify\_anchors.py}
\cite{YM04code}.

\begin{definition}\label{def:basis}
For $L\ge5$ let $|p\rangle:=\chi_{1/2}(U_p)$, $p\in\mathcal P$, the
orthonormal basis of the $E=3$ eigenspace $\mathcal M$; $P$ denotes the
orthogonal projection onto $\mathcal M$. For $L=4$ add winding states $|w\rangle$.
\end{definition}

\section{First order}\label{sec:first}
Write $V:=H_M-2|\mathcal P|\,\mathbf 1=-\sum_{q}\chi(U_q)$; $H(\mu)=H_E+2|\mathcal P|\mu+\mu V$.
\begin{lemma}[first-order cancellation]\label{lem:first}
$PVP=0$ on $\mathcal M$, and $\langle\Omega|V|\Omega\rangle=0$; moreover
$\langle p|V|\Omega\rangle=-\delta_{pq}$ summed, i.e.\
$\langle p|V|\Omega\rangle=-1$ for each $p$.
\end{lemma}
\begin{proof}
$\langle p|V|q\rangle=-\sum_r\int\chi(U_p)\,\chi(U_r)\,\overline{\chi(U_q)}\,dU$.
The integrand has total degree $(2,1)$ in the fundamental matrix coefficients
(one $\bar U$-type factor from $\overline{\chi(U_q)}$, two $U$-type).
Under the center element $z=-\mathbf 1\in\G$ the integrand acquires a factor
$z^{2}\bar z=z$ (i.e.\ $-1$), while the measure is invariant; hence every
integral vanishes. The same parity argument kills
$\langle\Omega|V|\Omega\rangle$ (degree $(1,0)$). For
$\langle p|V|\Omega\rangle=-\sum_r\int\chi(U_p)\overline{\chi(U_r)}\,dU=-\sum_r\delta_{pr}=-1$
by Lemma \ref{lem:gram}.
\end{proof}
Thus the first-order correction to the band is exactly $0$, and the first
correction to the ground energy is second order (Section~6).

\section{Second-order channels}\label{sec:channels}
Exhaustion of the intermediate spectrum --- that the four channels below are
\emph{all} the intermediate states with nonzero matrix elements --- is not
assumed; it is proved in Appendix~\ref{app:F},
Proposition~\ref{prop:exhaust}, from the Peter--Weyl selection rules and a
lattice-graph lemma.
The second-order effective operator on $\mathcal M$ is
\[
M_2 \;=\; P\,V\,Q\,(3-H_E)^{-1}Q\,V\,P ,
\]
where $Q=1-P$ and $(3-H_E)^{-1}$ acts on $Q\mathcal H_0$ (intermediate
states have $E_n\ge 4\cdot\tfrac12(\tfrac12+1)=3$, and $E_n=3$ states lie in
$\mathcal M$; the resolvent is negative). We classify the intermediate
states contributing to $\langle p|M_2|q\rangle$ by the support of the
inserted $\chi(U_r)$ relative to $p,q$.

\subsection{Vacuum intermediate}
$|\Omega\rangle\langle\Omega|$ contributes
$\frac{\langle p|V|\Omega\rangle\langle\Omega|V|q\rangle}{3-0}
=\frac{1}{3}$ for every pair $(p,q)$, i.e.\ the operator $+\frac13 J$ with
$J$ the all-ones matrix in the $|p\rangle$ basis.

\subsection{Single-plaquette ($\chi_1$) intermediate, $p=q$}\label{subsec:chi1}
With $r=p$ the insertion acts on $|p\rangle$ itself. Using the pointwise
fusion identity for $\G$ characters,
$\chi_{1/2}(g)^2=\chi_0(g)+\chi_1(g)$,
\[
-\chi(U_p)^2 \;=\; -\mathbf 1 \;-\; \chi_1(U_p),
\qquad \chi_1(U_p):=\Tr_{\mathrm{sym}^2}\bigl(U_{\ell_1}\cdots U_{\ell_4}\bigr).
\]
The constant term is the vacuum intermediate of Section 5.1. The state
$|\chi_1(p)\rangle:=\chi_1(U_p)$ is gauge-invariant, carries pure flux
$j=1$ on each of the four links of $p$ (each link appears once, in the
spin-$1$ representation), hence has $E_n=4\cdot 1\cdot 2=8$ and, by the
same Peter--Weyl chain as Lemma \ref{lem:gram}, norm $1$. Thus
$V|p\rangle\supset-|\chi_1(p)\rangle$, the regeneration amplitude is
$\langle p|V|\chi_1(p)\rangle=-1$, and the channel contributes
$(-1)^2/(3-8)=-\tfrac15$ on the diagonal.

\subsection{Touching pairs, $p\ne q$ sharing one link}\label{subsec:touch}
Two plaquettes sharing a link: inserting $\chi(U_r)$ with $r$ supported on
$p\cup q$ produces intermediate states on the $7$ links of $p\cup q$ in
which the shared link carries fused flux $s=0$ or $s=1$
($V_{1/2}\otimes V_{1/2}=V_0\oplus V_1$) and the other six links carry
$j=\tfrac12$. The two channels have energies
$E_s=6\cdot\tfrac34+s(s+1)$, i.e.\ $E_0=\tfrac92$, $E_1=\tfrac{13}{2}$
(denominators $-\tfrac32$, $-\tfrac72$), with squared amplitudes
$c_0^2=\tfrac14$, $c_1^2=\tfrac34$ (Appendix~B, derived twice: $6j$
recoupling and Weingarten; $c_0^2+c_1^2=1$ is the regeneration sum rule).
With regeneration amplitudes $-c_s$, the touching-pair matrix element is
\[
v\;=\;\frac{c_0^2}{3-E_0}+\frac{c_1^2}{3-E_1}
\;=\;-\frac{2}{3}\cdot\frac14-\frac{2}{7}\cdot\frac34
\;=\;-\frac16-\frac{3}{14}\;=\;-\frac{8}{21}.
\]
On the cubic lattice every bond lies in exactly four plaquettes (two planes
through the bond, two sides of each), so each plaquette shares one link with
$4\times 3=12$ others and no two distinct plaquettes share two links; the
touching adjacency $A$ is $12$-regular (machine census, Appendix~C), and the
touching channels contribute $v(A+12I)$.

\subsection{Disjoint pairs}\label{subsec:disjoint}
If $p,q$ are distinct and share no link, the intermediate state supported on
$p\cup q$ (regeneration, $E_n=6$, denominator $-3$) contributes $-\frac13$
off-diagonal. The diagonal disjoint bookkeeping contributes
$-\frac{|\mathcal P|-13}{3}$: the sum over $q$ disjoint from $p$ has
$|\mathcal P|-13$ terms ($12$ touching plus $p$ itself excluded).

\section{Effective Hamiltonian, vacuum subtraction, and the head of the band}
\label{sec:eff}
Summing the channels, with $A$ the touching adjacency and $J$ all-ones,
\begin{equation}\label{eq:M2}
M_2 \;=\; v(A+12I)\;+\;\tfrac13 J\;-\;\tfrac15 I\;-\;\tfrac13\bigl(J-I-A\bigr)
\;-\;\tfrac{|\mathcal P|-13}{3}\,I
\;=\;\Bigl(\tfrac13+v\Bigr)A+c\,I,
\end{equation}
$c=\tfrac13+12v-\tfrac15-\tfrac{|\mathcal P|-13}{3}$: the all-ones operator
cancels \emph{as an operator identity}. This identity is verified exactly at
$L=5$ ($|\mathcal P|=375$) in \cite{YM04code}, symbolically over the
plaquette link structure.

\subsection{Ground-state subtraction and the vacuum-mixing consistency}
The ground energy is, by non-degenerate analytic perturbation theory with
intermediate states $|p\rangle$ (amplitude $-1$, denominator $0-3$),
\[
E_0(\mu)=\mu^2\sum_{p\in\mathcal P}
\frac{|\langle p|V|\Omega\rangle|^2}{0-3}+O(\mu^4)
=-\frac{|\mathcal P|}{3}\,\mu^2+O(\mu^4) ;
\]
the third-order term vanishes identically: its Rayleigh--Schr\"odinger
expression is a sum of $\langle\Omega|V|n\rangle\langle n|V|m\rangle
\langle m|V|\Omega\rangle/((0-E_n)(0-E_m))$, and
$\langle\Omega|V|a\rangle\ne0$ forces $a$ to be a single-plaquette state:
in the Peter--Weyl flux basis, any flux link outside $\mathrm{supp}\,r$ leaves
an unpaired $\bar U$-factor and any $r$-link without matching flux leaves an
unpaired $U$-factor (Appendix~A, rule (ii)), so $\mathrm{supp}\,a=
\mathrm{supp}\,r$ with flux $\tfrac12$ throughout and $a\propto|p\rangle$ by
Lemma 3.2 (at $L=4$ a winding support never coincides with a plaquette
support); the middle factor is then some $\langle p|V|q\rangle=0$ (Lemma
\ref{lem:first}); hence the $O(\mu^4)$ bound. Equivalently, in the one-particle sector the uniform mode
$|u\rangle=|\mathcal P|^{-1/2}\sum_p|p\rangle$ mixes with the vacuum through
$\langle u|V|\Omega\rangle=-|\mathcal P|^{1/2}$; the $2\times2$ block
$\begin{pmatrix}0&-\mu|\mathcal P|^{1/2}\\ -\mu|\mathcal P|^{1/2}&3\end{pmatrix}$
has upper eigenvalue $3+\frac{|\mathcal P|}{3}\mu^2+O(\mu^4)$ and lower
eigenvalue $-\frac{|\mathcal P|}{3}\mu^2+O(\mu^4)$. The gap is
band-bottom minus $E_0(\mu)$; the $+\frac{|\mathcal P|}{3}$ entering the
$k=0$ band edge through $\frac13J$ in \eqref{eq:M2} is \emph{exactly} the
same vacuum-intermediate contribution, counted once in $M_2$ and cancelled
once by the subtraction $-E_0''=-\frac{|\mathcal P|}{3}$:
\begin{equation}\label{eq:c2}
c_2 \;=\; \lambda_{\min}\bigl((\tfrac13+v)A\bigr)+c+\frac{|\mathcal P|}{3}
\;=\;12\Bigl(\tfrac13+v\Bigr)+\tfrac13+12v-\tfrac15+\tfrac{13}{3}
\;=\;24v+\frac{127}{15}.
\end{equation}
Since $\tfrac13+v=-\tfrac1{21}<0$, the minimum of $(\tfrac13+v)A$ over the
band is $(\tfrac13+v)\lambda_{\max}(A)$, attained at $k=0$.

\begin{lemma}[band head]\label{lem:head}
For every $L\ge4$, $\lambda_{\max}(A(k))\le 12$ for all crystal momenta
$k\in[-\pi,\pi]^3$, with equality iff $k=0$; $\lambda_{\max}(A(0))=12$ with
the uniform orientation eigenvector.
\end{lemma}
\begin{proof}
Fourier transform in the translations: the fiber $A(k)$ is a $3\times3$
matrix indexed by plaquette orientation $(xy,xz,yz)$:
$A_{xy,xy}=2\cos k_x+2\cos k_y$, etc.,\ with off-diagonal
$A_{xy,xz}=4\cos(k_y/2)\cos(k_z/2)$ and cyclic. By Gershgorin, every
eigenvalue is $\le (2\cos k_x+2\cos k_y)+|A_{xy,xz}|+|A_{xy,yz}|
\le 4+8=12$, and equality forces $\cos k_x=\cos k_y=1$ and
$|\cos(k_y/2)|=|\cos(k_z/2)|=1$ etc.,\ i.e.\ $k=0$ in the Brillouin zone.
At $k=0$, $A(0)$ is the adjacency of the finite graph and the all-ones
vector gives $\lambda=12$ because $A$ is $12$-regular.
\end{proof}
The \emph{uniqueness} statements --- simplicity of the Perron root,
rigidity of equality $\lambda_{\max}(A(k))=12\Rightarrow k=0$, and the fact
that the exact deformed band minimum sits at $k=0$ for all sufficiently
small $\mu$, with a strictly positive quadratic gap of leading coefficient
$\tfrac{4}{63}$ to the next momentum --- are proved in
Appendix~\ref{app:E} (Perron--Frobenius plus exact curvature
$-\tfrac43|k|^2$; machine-checked in \texttt{code/verify\_bandhead.py}).

\begin{proof}[Proof of Theorem \ref{thm:main}]
$L\ge5$: Lemma \ref{lem:flux} gives $\mathcal M$ = plaquette states,
orthonormal (Lemma \ref{lem:gram}). First order vanishes
(Lemma \ref{lem:first}). Degenerate perturbation theory to second order
is regular because the cluster $E=3$ is separated from the rest of the
spectrum of $H_E$ on the gauge-invariant space by $\tfrac32$: the spectrum
of $H_E$ is $\sum_\ell j_\ell(j_\ell+1)$ over spin networks, and the vertex
rules (Lemma \ref{lem:valence}) plus bipartiteness of the cubic lattice
force every closed network with $E<\tfrac92$ other than the vacuum to be a
$4$-cycle at spin $\tfrac12$ ($2$-valent vertices force constant spin along
loops, odd loops do not exist, and the smallest theta graph carries
$E\ge7$); the next level is the $6$-loop $E=\tfrac92$. Thus the cluster is
isolated by a gap $\tfrac32$ uniformly in $L$, and the band functions are
$E_\alpha(\mu,k)=3+\mu^2\lambda_\alpha(M_2(k))+O(\mu^3)$ uniformly in $k$
(analyticity, Section 2.3). By \eqref{eq:M2}, \eqref{eq:c2} and
Lemma \ref{lem:head}, the bottom of the band minus $E_0(\mu)$ is
$3+\bigl(24v+\tfrac{127}{15}\bigr)\mu^2+O(\mu^3)$ with $v=-\tfrac8{21}$:
$24v+\tfrac{127}{15}=-\tfrac{192}{21}+\tfrac{127}{15}
=-\tfrac{64}{7}+\tfrac{127}{15}=\tfrac{-960+889}{105}=-\tfrac{71}{105}$.
The $|\mathcal P|$ dependence cancels identically in \eqref{eq:c2}; hence the
coefficient is $L$-independent for $L\ge5$.
\end{proof}

\section{Volume independence}\label{sec:volume}
The cancellation of $|\mathcal P|$ in \eqref{eq:c2} is exact, not asymptotic:
the extensive vacuum-intermediate term $\frac13|\mathcal P|$ on the uniform
mode is cancelled by the extensive diagonal disjoint term
$-\frac{|\mathcal P|-13}{3}$ plus the subtraction $+\frac{|\mathcal P|}{3}$.
What remains, $24v+\tfrac{127}{15}$, involves only the local touching
geometry, identical for all $L\ge5$ (the $13$ = $12+1$ touching census is
$L$-independent; at $L=4$ the winding sector changes the spectrum, Section 8).
The machine verification at $L=5$ builds the full $375\times375$ operator
over $\mathbb Q$ and confirms \eqref{eq:M2} entry-by-entry
\cite{YM04code}.

\section{The $L=4$ winding sector}\label{sec:L4}
At $L=4$ there are $192$ plaquettes and $48$ straight winding loops
($3$ directions $\times\ 4^2$ transverse positions); the census is exact and
machine-checked. Winding states $|w\rangle$ are orthonormal to plaquette
states and to each other by Lemma \ref{lem:gram}. First order remains zero
(center parity, Lemma \ref{lem:first}). \begin{lemma}[winding decoupling]\label{lem:mix}
$\langle p|M_2|w\rangle=0$ for every plaquette $p$ and winding loop $w$ at
$L=4$.
\end{lemma}
\begin{proof}
First the geometric census: $|p\cap w|\le1$ for all pairs --- the two
parallel links of a plaquette share their longitudinal coordinate, the four
links of a straight winding have distinct longitudinal coordinates (exact
census over all $192\times48$ pairs in \texttt{verify\_census.py}). Every
intermediate contributing to $\langle p|M_2|w\rangle$ arises from $V|w\rangle$:
the vacuum (killed by $\langle\Omega|V|w\rangle=0$, support: no plaquette
support equals a winding support), $|\chi_1(w)\rangle$ (energy $8$), fused
two-loop states on $w\cup r$ (energies $9/2,13/2$), or products $\chi_r\chi_w$
(energy $6$). Pairing with $\langle p|V$ inserts one further plaquette
character $\chi_s$, and covering all four winding links would require
$|p\cap w|+|s\cap w|\ge4$ while the census gives $\le2$; some winding link is
therefore left uncovered, carrying a single unpaired $U$-factor in the fused
and product channels (rule A.1(ii)) or two $U$-factors in one symmetric
spin-$1$ slot in the $\chi_1$ channel (Lemma \ref{lem:sym2}). Every channel
vanishes link by link.
\end{proof}
Thus the second-order block is diagonal in the plaquette/winding split. For winding states: $16$ touching channels (each
$=v$), no $\chi_1$ self channel, no vacuum channel
($\langle w|V|\Omega\rangle$ contributes only through the same parity
argument), disjoint diagonal $-\frac{|\mathcal P|-16}{3}$, giving
\[
c_{2,w}^{\rm diag}=16v+\frac{16}{3}=-\frac{16}{21}\approx-0.7619
<-\frac{71}{105}\approx-0.6762 .
\]
The winding chain block has zero diagonal, so the head of the winding band is
$\le-16/21$, strictly below $-71/105$: at $L=4$ the gap is governed by the
winding sector and Theorem \ref{thm:main} fails. The exact
$c_2(L{=}4)=-\frac{16}{21}+\lambda_{\min}(B_w)$ requires the two-link chain
recoupling of the block $B_w$, which we have not computed; we state the
bound and name the open value rather than guess it.

\section{Scope: fixed volume, uniformity, continuum}\label{sec:scope}
Theorem \ref{thm:main} is a fixed-$L$ perturbative identity:
$\mu_0$ may depend on $L$, and the $O_L(\mu^3)$ remainder is not controlled
uniformly in $L$. It does not imply $\Delta_{\rm unif}(\mu)>0$, nor anything
about $a\to0$. The distinction between (i) the finite-volume gap,
(ii) the volume-uniform gap at fixed coupling, and (iii) the continuum gap is
enforced throughout, following the campaign's model specification
\cite{YM01spec}. The relation of the coefficient to published
strong-coupling series (Kogut--Sinclair--Susskind and successors) and the
novelty classification are recorded in the literature ledger \cite{YM04ledger};
absence of an exact printed match is reported as ``not located,'' not as
proof of novelty. For honesty about the ingredients we note explicitly that
the individual second-order channel values $-\tfrac13$, $-\tfrac16$,
$-\tfrac{3}{14}$ appearing in Section~5 are published, in print, in the
closest fully accessible relative of our calculation --- the $2{+}1$-dimensional
strong-coupling static-potential analysis of Morningstar, whose staple
classes carry exactly these numbers (the recoupling weights
$\tfrac14,\tfrac34$ and denominators $-\tfrac32,-\tfrac72$ are
dimension-independent); what is not located in the accessible literature is
the assembled $3{+}1$-dimensional one-glueball-band coefficient
$-\tfrac{71}{105}$ with exact vacuum subtraction and the finite-volume
band-head theorem. The novelty classification is \textsc{Novelty
Unresolved}: the single object-matched prior source,
Kogut--Sinclair--Susskind, Nucl.\ Phys.\ \textbf{B114} (1976) 199--236,
remains \textsc{Access-Incomplete} after eleven independent retrieval
routes. The exact equivalence test is recorded for whoever obtains the
text: convert their printed second-order gap, under $\mu=2g^{-6}$ and
energies in units $g^2/2$, to
$\Delta_{\rm phys}=(3/2)g^{2}-(142/105)g^{-10}$; agreement yields
\textsc{Known}, a different converted value yields \textsc{Novelty
Supported}. No originality is claimed here.

\section*{Acknowledgments}
Prepared as mission YM-04 of a personal research campaign; the manuscript
supersedes nothing in the archive, which is preserved unchanged.

\begin{thebibliography}{99}
\bibitem{KS75} J.~B.~Kogut, L.~Susskind, ``Hamiltonian formulation of
lattice gauge theories,'' Phys.\ Rev.\ D \textbf{11} (1975) 395.
\bibitem{KSS76} J.~B.~Kogut, D.~K.~Sinclair, L.~Susskind, ``Monte Carlo
calculation of glueball energies in Hamiltonian lattice gauge theory,''
Nucl.\ Phys.\ \textbf{B114} (1976) 199.
\bibitem{DZ83} J.-M.~Drouffe, C.~Itzykson (and C.~B.~Lang), strong-coupling
methods; J.-M.~Drouffe, H.~J.~Zuber, Phys.\ Rep.\ \textbf{102} (1983) 1.
\bibitem{MM94} C.~Montvay, G.~M\"unster, \emph{Quantum Fields on a
Lattice}, Cambridge (1994).
\bibitem{Weyl} H.~Weyl, \emph{The Classical Groups}; Peter--Weyl
orthogonality as in M.~R\"oser, or standard references; Weingarten calculus:
D.~Weingarten, J.\ Math.\ Phys.\ \textbf{19} (1978) 999; Collins-type bounds
not needed (exact finite computations only).
\bibitem{YM04code} Verification code (exact-rational SymPy, no numerical
fitting): \texttt{code/verify\_anchors.py}, \texttt{verify\_identity.py},
\texttt{verify\_census.py}, \texttt{verify\_bandhead.py}, accompanying this
manuscript.
\bibitem{YM04ledger} Supplementary documents accompanying this manuscript
(static copies, checksummed in \texttt{MANIFEST.sha256}): the primary-source
access log for \textsc{B114} (1976) with the recorded equivalence test, and
the gap taxonomy (fixed-volume / volume-uniform / continuum) used in
Section~9. The scope and novelty statements of Section~9 are reproduced in
full in the body of this paper, so the manuscript is logically
self-contained.
\end{thebibliography}

\appendix
\section{Peter--Weyl index chains and vanishing lemmas}\label{app:A}
Write $\chi(g)=\Tr g$ (spin $\tfrac12$ character). For a loop
$C=(\ell_1,\dots,\ell_4)$,
$\chi(U_C)=U_1{}^{a}{}_{b}\,U_2{}^{b}{}_{c}\,U_3{}^{c}{}_{d}\,
U_4{}^{d}{}_{a}$ (summation over repeated indices).

\begin{lemma}[vanishing rules]\label{lem:degree}
Let $F$ be a monomial in link matrix coefficients with per-link bidegree
$(p_\ell,q_\ell)$ (numbers of $U$-type and $\bar U$-type factors of link
$\ell$). Then $\int F\,dU=0$ unless
\begin{enumerate}
\item[(i)] $p_\ell-q_\ell$ is \emph{even} for every $\ell$ (single-link
center parity: the flip $U_\ell\mapsto -U_\ell$ preserves Haar measure and
multiplies $F$ by $(-1)^{p_\ell-q_\ell}$); and
\item[(ii)] the representation tensor product carries an invariant vector
after all intertwiners are contracted; in particular a link of bidegree
$(1,0)$ or $(0,1)$ forces vanishing, since $V_{1/2}$ has no invariant and
$\int U^a{}_b\,dU=0$.
\end{enumerate}
These are the only rules used in this paper. (Neither is a complete
criterion: unequal-degree integrals such as $\int UU$ \emph{can} be nonzero
for $\G$, via $\varepsilon$-pairings --- see Lemma \ref{lem:sym2}.)
\end{lemma}

\begin{lemma}[symmetric-slot rule]\label{lem:sym2}
$\displaystyle\int U^a{}_b\,U^c{}_d\,dU=\tfrac12\,\varepsilon^{ac}
\varepsilon_{bd}$; consequently, if a monomial has some link contributing
two $U$-factors occupying a single symmetric (spin-$1$) matrix coefficient
and no $\bar U$-factor on that link, the integral vanishes, because
$\varepsilon^{ac}$ annihilates the symmetric pair.
\end{lemma}
\begin{proof}
The left side is an intertwiner in $\mathrm{Inv}(V\otimes V\otimes V\otimes V)$
unique up to scale, and $\tfrac12\varepsilon^{ac}\varepsilon_{bd}$ is the
invariant pairing; the annihilation statement is $\varepsilon^{ac}S_{ac}=0$
for symmetric $S$.
\end{proof}

\begin{lemma}[diagonal Gram]\label{lem:gramA}
$\langle\chi(U_C),\chi(U_C)\rangle=1$.
\end{lemma}
\begin{proof}
$\overline{\chi(U_C)}=\chi(U_C^\dagger)
=\bar U_4{}^{a}{}_{d}\bar U_3{}^{d}{}_{c}\bar U_2{}^{c}{}_{b}\bar U_1{}^{b}{}_{a}$.
Apply Peter--Weyl orthogonality link by link,
$\int U^a{}_b\,\overline{U^c{}_d}\,dU=\tfrac12\delta^a_d\,\delta^c_b$:
each of the four links contributes $\tfrac12$; each of the four successive
index closures contributes a factor $2$; the product is
$2^{-4}\cdot 2^{4}=1$.
\end{proof}

\begin{lemma}[off-support vanishing]\label{lem:offsupport}
If $C\ne C'$ are closed loops of length $4$ on $\mathbb T_L^3$, then
$\langle\chi(U_C),\chi(U_{C'})\rangle=0$.
\end{lemma}
\begin{proof}
Some $\ell\in C\setminus C'$; the integrand has bidegree $(1,0)$ on $\ell$,
so Lemma \ref{lem:degree}(ii) applies.
\end{proof}

\section{Recoupling coefficients $c_0^2=\tfrac14$, $c_1^2=\tfrac34$}\label{app:B}
\subsection{Method 1: $6j$/F-move calculus}
Let $|s\rangle$, $s\in\{0,1\}$, be the orthonormal $s$-channel basis of the
intertwiner space $\mathrm{Inv}(V_{1/2}^{\otimes 4})$: pair links
$(\ell_1,\ell_2)$ to spin $s$ and $(\ell_3,\ell_4)$ to the same $s$. The
regeneration state $\chi(U_{p})\,\chi(U_{q})$ for touching $p,q$ (shared
link $\ell_1$) is, in link language, the four-valent contraction with
$\ell_1$ appearing twice; changing from the plaquette pairing to the
shared-link pairing is the $F$-move
\[
F=\begin{pmatrix}\tfrac12 & \tfrac{\sqrt3}{2}\\[2pt]
\tfrac{\sqrt3}{2} & -\tfrac12\end{pmatrix},
\qquad
F_{ss'}=\sqrt{(2s+1)(2s'+1)}
\left\{\begin{matrix}\tfrac12&\tfrac12&s\\ \tfrac12&\tfrac12&s'\end{matrix}\right\},
\]
i.e.\ the $6j$ values
$\{ \tfrac12\tfrac12 0;\tfrac12\tfrac12 0\}=\tfrac12$,
$\{\tfrac12\tfrac12 1;\tfrac12\tfrac12 1\}=-\tfrac16$,
$\{\tfrac12\tfrac12 0;\tfrac12\tfrac12 1\}=\tfrac12$
($\pm$ signs convention-dependent; only squares enter). Hence the
touching-regeneration state decomposes with amplitudes
$c_0=F_{00}=\tfrac12$, $c_1=F_{10}=\tfrac{\sqrt3}{2}$:
$c_0^2=\tfrac14$, $c_1^2=\tfrac34$, $c_0^2+c_1^2=1$ (unit regeneration norm,
Lemma \ref{lem:gramA}). Energies: the $s$-channel intermediate has six links
of flux $\tfrac12$ and the shared link of flux $s$:
$E_s=6\cdot\tfrac34+s(s+1)\in\{\tfrac92,\tfrac{13}{2}\}$, so
$v=\frac{c_0^2}{3-9/2}+\frac{c_1^2}{3-13/2}=-\frac8{21}$.

\subsection{Method 2: Weingarten integration}
$\langle p|V\,\Pi_s\,V|q\rangle$ for touching $p,q$ reduces to integrals of
products of four matrix coefficients on the shared link,
$\int U^a{}_bU^c{}_d\bar U^e{}_f\bar U^g{}_h\,dU$, evaluated by the $\G$
Weingarten rule
$\int U\otimes\bar U\otimes U\otimes\bar U
=\tfrac{1}{6}(\text{id}-\tfrac12\text{flip})
+\tfrac{1}{3}(\text{flip}-\tfrac12\text{id})$
(dimension formulas $\int P_{\mathrm{sym}}=\tfrac34$,
$\int P_{\mathrm{asym}}=\tfrac14$). Both methods agree exactly; the
machine check is \texttt{verify\_anchors.py} and
\texttt{verify\_identity.py} \cite{YM04code}.

\subsection{$\chi_1$ regeneration}
Pointwise fusion $\chi_{1/2}(g)^2=\chi_0(g)+\chi_1(g)$ gives
$V|p\rangle\supset-\mathbf 1-|\chi_1(p)\rangle$; $|\chi_1(p)\rangle$ carries
pure flux $j=1$ on the four links of $p$ ($E_n=8$) and has norm $1$ by the
index chain of Lemma \ref{lem:gramA} with the spin-$1$ representation
($\int\chi_1\bar\chi_1$-type contraction, same $2^{\pm4}$ bookkeeping).
Amplitude $\langle p|V|\chi_1(p)\rangle=-1$, channel $-1/5$.

\section{Adjacency census and the band-head lemma}\label{app:C}
\textbf{Census.} On $\mathbb T_L^3$ every bond lies in exactly four
plaquettes (two planes through the bond direction, two transverse sides),
hence each plaquette shares exactly one link with $4\text{ (links)}\times
3\text{ (partners)}=12$ others, and distinct plaquettes never share two
links. The touching adjacency $A$ is $12$-regular; both the degree and the
$\{0,1\}$ intersection property are verified exhaustively at $L=5$
($n=375$) in \texttt{verify\_identity.py}.

\textbf{Fiber decomposition.} Translations act freely on plaquettes of each
of the three orientations $(xy,xz,yz)$; Fourier transform along
$\mathbb Z_L^3$ block-diagonalizes $A$ into $3\times3$ fibers
\[
A(k)=\begin{pmatrix}
2\cos k_x+2\cos k_y & 4\cos\tfrac{k_y}{2}\cos\tfrac{k_z}{2} &
4\cos\tfrac{k_x}{2}\cos\tfrac{k_z}{2}\\
4\cos\tfrac{k_y}{2}\cos\tfrac{k_z}{2} & 2\cos k_x+2\cos k_z &
4\cos\tfrac{k_x}{2}\cos\tfrac{k_y}{2}\\
4\cos\tfrac{k_x}{2}\cos\tfrac{k_z}{2} & 4\cos\tfrac{k_x}{2}\cos\tfrac{k_y}{2} &
2\cos k_y+2\cos k_z
\end{pmatrix},
\]
(derivation: the two bonds through which an $xy$-plaquette at the origin
touches $xz$-plaquettes are displaced by $\pm e_y$ and $\pm e_z$ relative to
the shared-bond identification, producing the product-of-cosines entries;
each entry is the sum of the four phase weights of the corresponding
touching family).

\textbf{Gershgorin bound.} Row $xy$:
$\lambda\le 2\cos k_x+2\cos k_y+4|\cos\tfrac{k_y}{2}\cos\tfrac{k_z}{2}|
+4|\cos\tfrac{k_x}{2}\cos\tfrac{k_z}{2}|\le 4+8=12$; equality forces
$\cos k_x=\cos k_y=1$ and $|\cos\tfrac{k_{y,z}}{2}|=1$, i.e.\ $k=0$ in
$[-\pi,\pi]^3$, where $A(0)\mathbf 1=12\,\mathbf 1$ by regularity. This is
Lemma \ref{lem:head} records the \emph{upper bound}; uniqueness of the
maximizer and of the exact band head are proved in Appendix~\ref{app:E},
without numerics.

\section{Normalization dictionary and reproducibility}\label{app:D}
\textbf{Dictionary.} The theorems are statements about
$H(\mu)=\sum_\ell C_2(\ell)+\mu\sum_p(2-\Re\Tr U_p)$.
For $H_{\rm KS}=\frac{g^2}{2}\sum E^2-\frac{1}{2g^4}\sum_p(\Tr U_p+\Tr
U_p^\dagger)$ at $a=1$: $\mu=2g^{-6}$, physical energy
$=\frac{g^2}{2}\times$(dimensionless). If an alternative convention
rescales the expansion parameter by a constant, $\mu'=c\mu$, the theorem
reads $\Delta=3-\frac{71}{105\,c^{2}}\mu'^{2}+O_L(\mu'^{3})$. Conventions
quoting $\mu'=g^{-4}$ arise from applying the same dimensionless form to
$\widetilde H=\frac{2}{g^2}H_{\rm KS}$ with $g$ reparametrized so that the
plaquette coefficient reads $-\mu'\sum_p\Re\Tr U_p$ with $\mu'=g^{-4}$
(a different overall normalization of $H_{\rm KS}$); the dictionary must then be
applied to the full dimensionless Hamiltonian, not to the symbol $\mu$
alone --- convert the Hamiltonian, not just the coupling.

\textbf{Reproducibility.} \texttt{code/verify\_anchors.py} proves
$\|\chi(U_p)\|^2=1$, $\langle p|V|\Omega\rangle=-1$,
$\langle p|V|q\rangle=0$, and the $c_s^2$ by two independent exact methods
(Weyl integration and Weingarten); \texttt{code/verify\_identity.py}
builds the $375\times375$ operator at $L=5$ over $\mathbb Q$, verifies
$M_2=(\tfrac13+v)A+cI$ entry-by-entry, the $2\times2$ mixing consistency
$\langle u|\tfrac13J|u\rangle=|\mathcal P|/3$, and $c_2=-71/105$;
\texttt{code/verify\_census.py} runs the $L=4$ census
$192+48=240$ and the winding touching count. All computations are exact
rational arithmetic (SymPy); no numerical fitting, no floating-point
assumptions.

\section{Band-head uniqueness: Perron--Frobenius and curvature}
\label{app:E}
\begin{lemma}[maximal eigenvalue, strict form]\label{lem:headE}
For every $L\ge4$ and every $k\in[-\pi,\pi)^3$,
$\lambda_{\max}(A(k))\le12$, with equality \emph{if and only if} $k=0$;
moreover $\lambda_{\max}(A(0))=12$ is \emph{simple} and its eigenline is
spanned by the uniform orientation vector.
\end{lemma}
\begin{proof}
\emph{(1) Comparison.} For a Hermitian matrix $B$, writing $|B|$ for the
entrywise absolute value,
$\lambda_{\max}(B)=\max_{\|x\|=1}\langle x,Bx\rangle\le
\max_{\|x\|=1}|\langle x,Bx\rangle|\le
\max_{\|y\|=1,\,y\ge0}\langle y,|B|y\rangle=\rho(|B|)$,
the last step being the Collatz--Wielandt characterization of the Perron
root of a nonnegative matrix. Applied to $B=A(k)$: every entry of $|A(k)|$
is bounded by the corresponding entry of $A(0)$ (diagonals:
|$2\cos k_x+2\cos k_y|\le4$; off-diagonals: products of cosines bounded by
$4$), and the spectral radius is monotone on nonnegative matrices, so
$\lambda_{\max}(A(k))\le\rho(|A(k)|)\le\rho(A(0))$.

\emph{(2) Irreducibility.} The touching graph of plaquettes on $\mathbb T_L^3$
is connected: in one plane the plaquettes form the $L\times L$ grid under
link-sharing, and the three plane-families meet at every site. Hence $A(0)$
is irreducible; $12$-regularity gives $A(0)\mathbf1=12\mathbf1$, so
$\rho(A(0))=12$, and Perron--Frobenius makes $12$ a simple eigenvalue with
strictly positive eigenvector; regularity then forces that eigenvector to be
$\mathbf1$.

\emph{(3) Rigidity of equality.} Suppose $\lambda_{\max}(A(k))=12$. Then
equality holds throughout (1), in particular $\rho(|A(k)|)=\rho(A(0))$. If
some entry of $|A(k)|$ were strictly smaller than that of $A(0)$, then,
$A(0)$ being irreducible, the strict monotonicity of the Perron root (add a
nonnegative perturbation supported on that entry and use the Perron eigenvector
of $|A(k)|$, which is positive on the irreducible support) gives
$\rho(|A(k)|)<\rho(A(0))$, contradiction. Hence $|A(k)|=A(0)$ entrywise.
The off-diagonal equalities give $|\cos\tfrac{k_i}{2}|=1$ for each $i$;
for $k_i\in[-\pi,\pi)$ this forces $k_i=0$ (the diagonal equalities then
hold automatically).
\end{proof}

\begin{proposition}[exact band head for small $\mu$]\label{prop:headexact}
Let $L\ge5$. There are $\mu_1(L)>0$ and $\kappa(L)>0$ such that for
$0<\mu<\mu_1(L)$ the minimum of the deformed $E=3$ band over the discrete
momentum lattice $k\in\frac{2\pi}{L}\mathbb Z^3\cap[-\pi,\pi)^3$ is attained
uniquely at $k=0$, in the uniform orientation mode, and the gap to the next
band minimum is at least $\kappa(L)\mu^2$.
\end{proposition}
\begin{proof}
By Lemma \ref{lem:headE}(2) the eigenvalue $12$ of $A(0)$ is simple, so
$\lambda_{\max}(A(k))$ is real-analytic near $k=0$ (analytic perturbation of
a simple eigenvalue of the $3\times3$ fiber). Differentiating the
characteristic polynomial $p(k,\lambda)=\det(\lambda-A(k))$ at $(0,12)$:
with $A(k)=A(0)+A_2(k)+O(|k|^4)$, $A_2$ collecting the quadratic terms
($\operatorname{diag}$ entries $-k_x^2-k_y^2$ etc.,\ off-diagonal entries
$-\tfrac12(k_y^2+k_z^2)$ etc.), first-order perturbation theory on the simple
 eigenpair $(12,\mathbf1/\sqrt3)$ gives
$\lambda_{\max}(A(k))=12+\langle\mathbf1,A_2(k)\mathbf1\rangle/3+O(|k|^4)
=12-\tfrac{4}{3}|k|^2+O(|k|^4)$: the diagonal part contributes
$-\tfrac13\cdot2|k|^2$ and the off-diagonal part $-\tfrac13\cdot2|k|^2$,
summing to the isotropic $-\tfrac43|k|^2$ (exact rational; the cubic term
vanishes by the inversion symmetry $k\mapsto-k$ of the fiber; verified in
\texttt{code/verify\_bandhead.py}).
Since $\tfrac13+v=-\tfrac1{21}<0$, the lowest band function obeys,
uniformly in $k$ (Analyticity subsection of Sec.~\ref{sec:setup}),
$E_{\min}(\mu,k)=3+\mu^2[(\tfrac13+v)\lambda_{\max}(A(k))+c]+O(\mu^3)$ with
the $O(\mu^3)$ constant depending only on $L$ (finite lattice, bounded
perturbation). Hence for $k\ne0$,
$E_{\min}(\mu,k)-E_{\min}(\mu,0)
=\mu^2\cdot\tfrac{1}{21}\cdot\tfrac{4}{3}|k|^2+O(\mu^2|k|^4)+O(\mu^3)
=\tfrac{4}{63}\mu^2|k|^2+O(\mu^2|k|^4)+O(\mu^3)$.
On the discrete lattice the smallest nonzero $|k|$ is $2\pi/L$, so choosing
$\mu_1(L)$ with $C_{\rm rem}\,\mu_1<\tfrac{4}{63}(2\pi/L)^2$ makes the right
side strictly positive for every $k\ne0$, with margin $\kappa(L)\mu^2$.
Uniqueness of the maximizer at $k=0$ is Lemma \ref{lem:headE}(2).
\end{proof}

\section{Exhaustion of second-order intermediate states}\label{app:F}
The second-order effective operator on the $E=3$ eigenspace is
$M_2=-P\,H_M\,R_0\,H_M\,P$ restricted to the cluster with the vacuum
channel separated, $R_0$ the reduced resolvent of $H_E$ at $3$ on
$\mathcal M\ominus$ vacuum directions; every claim about
$M_2$ below is about matrix elements of this reduced operator between
plaquette states. The intermediate-state enumeration in
Proposition~\ref{prop:exhaust} is stated for the \emph{unreduced}
second-order expression $P\,H_M\,(H_E-3)^{-1}H_M\,P$ before vacuum
subtraction; channel (i) there is precisely the term removed by the
subtraction and reappears as the normalization correction of
Section~5.1.

\begin{lemma}[intermediate selection]\label{lem:exhaust1}
If $\langle n|H_M|p\rangle\ne0$ for a plaquette state $|p\rangle$ and a
gauge-invariant eigenstate $|n\rangle$ of $H_E$, then for some plaquette
$q$, $|n\rangle$ has an index chain supported in
$\mathrm{supp}(p)\cup\mathrm{supp}(q)$ in which every link of
$\mathrm{supp}(q)\setminus\mathrm{supp}(p)$ carries spin $\tfrac12$, every
link of $\mathrm{supp}(p)\cap\mathrm{supp}(q)$ carries spin $0$ or $1$, and
every link of $\mathrm{supp}(p)\setminus\mathrm{supp}(q)$ carries spin
$\tfrac12$. In particular $E(n)=\sum j(j{+}1)$ over that support.
\end{lemma}
\begin{proof}
$\langle n|H_M|p\rangle=\sum_{q}\int\overline{\chi_1(n)}\,\chi_{1/2}(U_q)\,
\chi_1(p)\,dU$. For each link $\ell$, the two plaquette
characters $\chi_1(p)\,\chi_{1/2}(U_q)$ carry, in their tensor-network
expansion, one fundamental slot per link around the closed chain of $p$
and one around the chain of $q$.
\footnote{\textbf{Erratum (2026-11-06).} The round-3 text read ``two
fundamental slots around the closed spin-$1$ chain of $p$'' and counted
``three slots (odd)'' on the overlap. Under the frozen convention
$|p\rangle=\chi_{1/2}(U_p)$ (Sec.~2; the index display of Appendix~A has
one slot per link) each character contributes \emph{one} slot per link,
so the overlap carries \emph{two} slots. The selection-rule
\emph{statement} is unaffected; the standalone derivation in
Appendix~G uses the corrected count. Correction history preserved in
\texttt{evidence/REFEREE\_R3PRIME\_ADJUDICATION.md}.} A link in
$\mathrm{supp}(q)\setminus\mathrm{supp}(p)$ carries exactly one
fundamental pair; degree matching of the $U_\ell$ integral against the
conjugate spin-$j_\ell(n)$ chain of $\overline{\chi_1(n)}$ then forces
$j_\ell(n)=\tfrac12$. A link in $\mathrm{supp}(p)\cap\mathrm{supp}(q)$
carries two slots; the admissible intermediate spins there are the
recoupling channels of the two fundamental legs of $q$ against the
spin-$1$ chain of $p$ computed in Appendices A--B: the singlet channel
($s=0$, weight $c_0^2=\tfrac14$) and the triplet channel ($s=1$, weight
$c_1^2=\tfrac34$), so $j_\ell(n)\in\{0,1\}$ on the overlap.
\textbf{Resolved (revision item R3'):} the exhaustive form of this
selection rule --- that the recoupling channels of Appendices A--B are
the \emph{only} nonvanishing overlap contributions, with no higher-spin
chain terms --- is now established by the standalone Peter--Weyl
derivation of Appendix~G (internally proved; independently
model-reviewed, not human peer review). The prior text read: ``\ldots\ is
currently established by the exact full-space operator identity of
Appendix D and the complete vertex-rule enumeration of
Proposition~\ref{prop:exhaust}, not by a standalone link-level parity
argument; a self-contained derivation is required before external review
of the exhaustion proof.'' Rule (ii) then requires an intertwiner
at each site of the union support; sites outside
$\mathrm{supp}(p)\cup\mathrm{supp}(q)$ carry unmatched single $\bar U$
factors from $\overline{\chi_1(n)}$ unless $j_\ell(n)=0$ there, giving the
support claim.
\end{proof}

\begin{lemma}[geometric trichotomy]\label{lem:exhaust2}
For plaquettes $p\ne q$ on $\mathbb T_L^3$, $|\mathrm{supp}(p)\cap
\mathrm{supp}(q)|\in\{0,1\}$: distinct plaquettes share no link or exactly
one link, never two. Hence the pairs $(p,q)$ contributing to
$\langle p|M_2|p'\rangle$ fall into exactly three geometric classes:
$q=p$; $q$ touching $p$ (one shared link); $q$ disjoint from $p$.
\end{lemma}
\begin{proof}
Two distinct unit squares of the cubic lattice sharing two edges would
share two adjacent edges, forcing a common vertex and two common edge
directions, hence a common plane and a common third and fourth vertex ---
i.e.\ the same square; or share two opposite edges, forcing the same four
vertices and the same plane. Both give $p=q$. On $\mathbb T_L^3$ with
$L\ge4$ the projection is injective on unit squares except for the winding
loops of Lemma (winding census), which are excluded from the plaquette set
by definition ($4<L$ edges). Exhaustive verification at $L=5$
($375$ plaquettes, all pairs) is in \texttt{verify\_identity.py}.
\end{proof}

\begin{lemma}[vertex singlet rules]\label{lem:valence}
For $j_1,j_2\in\tfrac12\mathbb N$: $\dim\mathrm{Inv}(j_1\otimes j_2)=1$ if
$j_1=j_2$ and $0$ otherwise. For three spins,
$\dim\mathrm{Inv}(j_1\otimes j_2\otimes j_3)=1$ if
$|j_1-j_2|\le j_3\le j_1+j_2$ and $j_1+j_2+j_3\in\mathbb Z$, and $0$
otherwise.
\end{lemma}
\begin{proof}
Clebsch--Gordan: $j_1\otimes j_2\simeq\bigoplus_{j=|j_1-j_2|}^{j_1+j_2}j$
(multiplicity-free); $\mathrm{Inv}$ picks the spin-$0$ summand. Two spins:
$0$ appears iff $j_1=j_2$. Three: fuse $j_1\otimes j_2$ and pick $j_3$ in
the range; the parity condition is $j_1+j_2+j_3\in\mathbb Z$.
\end{proof}

\begin{proposition}[channel exhaustion]\label{prop:exhaust}
The intermediate states with nonzero contribution to the unreduced
second-order sum before vacuum subtraction are exhausted by exactly the
four channels computed in Sections 5.1--5.4. Channel (i) is carried
entirely by the normalization (vacuum-subtraction) term of Section 5.1
and does not enter $M_2$; channels (ii)--(iv) exhaust the intermediates
of the reduced operator $\langle p|M_2|p'\rangle$:
\textup{(i)} the vacuum intermediate, value $-\tfrac13$ per $p$
(normalization only);
\textup{(ii)} the same-plaquette spin-$1$ regeneration, $E(n)=8$, value
$-\tfrac15$;
\textup{(iii)} touching pairs, two intermediate families with shared-link
spin $0$ ($E=\tfrac92$) and spin $1$ ($E=\tfrac{13}{2}$), each with a
unique intertwiner, combining through the Appendix-B weights
$c_0^2=\tfrac14$, $c_1^2=\tfrac34$ into $v=-\tfrac8{21}$;
\textup{(iv)} disjoint pairs, unique intermediate with all eight links at
spin $\tfrac12$, $E=6$, value $-\tfrac13$.
Every admissible intermediate below has intertwiner multiplicity exactly
$1$, so no multiplicity factors are hidden.
\end{proposition}
\begin{proof}
By Lemma \ref{lem:exhaust1} any intermediate chain is supported on
$\mathrm{supp}(p)\cup\mathrm{supp}(q)$ with spins in $\{0,1\}$ on
$\mathrm{supp}(p)\cap\mathrm{supp}(q)$ and $\tfrac12$ elsewhere; by
Lemma \ref{lem:exhaust2} the pair geometry is one of three cases, and the
admissible assignments are exactly those whose every vertex passes
Lemma \ref{lem:valence}. This is a finite check and it is executed exactly
(integer arithmetic) in \texttt{code/verify\_intertwiners.py}:
\emph{self} ($q=p$): $16$ assignments on the $4$-cycle; the four vertices
are $2$-valent in the support, so Lemma \ref{lem:valence} forces adjacent
spins equal, and going around the cycle leaves exactly two admissible
assignments --- all-$0$ ($E=0$, channel (i)) and all-$1$ ($E=8$, channel
(ii)); every mixed $0/1$ assignment fails at a vertex of unequal adjacent
spins.
\emph{Touching}: the shared link carries $s\in\{0,1\}$; the six singly
covered links carry $\tfrac12$; the four non-shared vertices are $2$-valent
with equal spins (pass), and the two endpoints of the shared link are
$3$-valent with $(s,\tfrac12,\tfrac12)$, admissible for both $s=0$ and
$s=1$ with multiplicity one, giving $E=\tfrac92$ and $E=\tfrac{13}{2}$
with denominators $-\tfrac32$ and $-\tfrac72$. The shared link enters
through the two singlet channels of the two fundamental legs at each
endpoint with the squared recoupling weights $c_0^2=\tfrac14$,
$c_1^2=\tfrac34$ of Appendix~B, so the summed touching value is
$v=\tfrac{1/4}{-3/2}+\tfrac{3/4}{-7/2}=-\tfrac16-\tfrac{3}{14}
=-\tfrac{8}{21}$, matching Section 5.3 exactly (machine-checked).
\emph{Disjoint}: both squares regenerate the all-$\tfrac12$ network
(unique invariant per square by the $2$-valent rule), $E=6$, denominator
$-3$, value $-\tfrac13$, matching Section 5.4 (machine-checked).
Multiplicities: $1$ (self), $12$ (touching, $12$-regularity),
$|\mathcal P|-13$ (disjoint), matching \eqref{eq:M2}. No other pair
geometry exists (Lemma \ref{lem:exhaust2}) and no other spin assignment on
these supports passes Lemma \ref{lem:valence}, so the list is exhaustive.
\end{proof}

\section{Standalone exhaustion proof of the selection rule (R3')}
\label{app:G}

This appendix supplies the self-contained derivation required by
revision item R3': the selection rule of Lemma~\ref{lem:exhaust1} is
exhaustive --- the singlet and triplet recoupling channels of
Appendices~A--B are the \emph{only} nonvanishing overlap contributions,
with no higher-spin chain terms. The argument uses only the frozen
conventions of Sec.~2, Schur--Weyl theory, and Peter--Weyl
orthogonality; it is independent of the full-space operator identity of
Appendix~D and of the enumeration of Proposition~\ref{prop:exhaust},
which now serve as cross-checks.

\begin{lemma}[External type of the intermediate states]
\label{lem:extype}
Let $\psi_n$ be a gauge-invariant $H_E$-eigenstate with per-link spins
$\{j_\ell(n)\}$. Regarded as a function on $G^L$, $G=\mathrm{SU}(2)$,
$\psi_n$ lies in the coefficient space of $\bigotimes_\ell
\rho_{j_\ell}$, $\rho_j=\mathrm{Sym}^{2j}$: under left $\times$ right
translation on link $\ell$ it is of \emph{external type} $(j_\ell,j_\ell)$
(diagonal); the vertex intertwiners supply the endpoint contractions.
Conjugation preserves external type ($V^*\cong V$ via $\varepsilon$), so
$\overline{\psi_n}$ has the same diagonal type, and reversed loop
traversal (whole-loop inversion, $\Tr A^{-1}=\Tr A$ on $\mathrm{SU}(2)$)
maps the slot reading $(2,0)\leftrightarrow(1,1)$, $(1,0)\leftrightarrow
(0,1)$ without changing the integrand or the admissible spins.
\end{lemma}

\begin{lemma}[Four-channel decomposition of the insertion]
\label{lem:fourt}
Decompose $F:=\chi_{1/2}(U_q)\,\chi_{1/2}(U_p)$ link by link in external
typing. A link in $\mathrm{supp}(p)\,\Delta\,\mathrm{supp}(q)$ carries a
single matrix entry, of type $(\tfrac12,\tfrac12)$. An overlap link
carries a product of two entries, content $(V\otimes V,\,
V^*\otimes V^*)$, and by Schur--Weyl $V\otimes V=0\oplus 1$ in each slot
independently:
\[
(0\oplus 1,\;0\oplus 1)\;=\;(0,0)\oplus(0,1)\oplus(1,0)\oplus(1,1),
\]
four channels, the crossed ones included. A link outside
$\mathrm{supp}(p)\cup\mathrm{supp}(q)$ contributes type $(0,0)$. The
decomposition is complete and multiplicity-free.
\end{lemma}

\begin{lemma}[Crossed-channel annihilation]
\label{lem:cross}
The Haar integral on $G^L$ is bi-invariant: applied to a pure-type tensor
it is nonzero only on the $(0,0)$ component of every link. Since
$\overline{\psi_n}$ has type $(j_\ell,j_\ell)$, the product
$\overline{\psi_n}\cdot F$ contains $(0,0)$ on link $\ell$ only where
$F$ has a component of type $(j_\ell,j_\ell)$: the crossed channels
$(1,0)$ and $(0,1)$ pair with $\overline{\psi_n}$ to $(j\otimes 1,\,
j\otimes 0)$ resp.~$(j\otimes 0,\, j\otimes 1)$, which contain $(0,0)$
only if $j=1$ and $j=0$ simultaneously --- impossible --- hence crossed
channels contribute zero for every $j$.
\end{lemma}

\begin{theorem}[R3', exhaustive form]
\label{thm:r3prime}
If $\langle n|H_M|p\rangle\ne 0$ then $j_\ell(n)=\tfrac12$ on
$\mathrm{supp}(q)\,\Delta\,\mathrm{supp}(p)$, $j_\ell(n)\in\{0,1\}$ on
$\mathrm{supp}(p)\cap\mathrm{supp}(q)$, and $j_\ell(n)=0$ outside
$\mathrm{supp}(p)\cup\mathrm{supp}(q)$; and the only nonvanishing
overlap contributions are the singlet and triplet channels of
Appendix~B, with weights $c_0^2=\tfrac14$, $c_1^2=\tfrac34$.
\end{theorem}
\begin{proof}
Necessity and exhaustivity: Lemmas~\ref{lem:extype}--\ref{lem:cross}
link by link --- symmetric-difference links force $j=\tfrac12$ (only
channel $(\tfrac12,\tfrac12)$); overlap links force $j\in\{0,1\}$
(diagonal channels $(0,0)$ and $(1,1)$; crossed channels annihilate);
exterior links force $j=0$. No other channels exist by the completeness
of the Schur--Weyl decomposition, so no higher-spin chain terms occur.
Realizability: Appendix~B computes both surviving projectors explicitly
($c_0^2=\tfrac14$, $c_1^2=\tfrac34$, both nonzero; $P_0+P_1=1$ on
$V\otimes V$), and the vertex rules (Lemma~\ref{lem:valence}) give
multiplicity one, reproducing the assignment list of
Proposition~\ref{prop:exhaust}: self $\{E=0,\;E=8\}$, touching
$\{E=\tfrac92,\;E=\tfrac{13}{2}\}$, disjoint $\{E=6\}$. All weights,
energies, content tables, and the symmetric-block $\varepsilon$-annihilation
are additionally machine-checked in exact rational arithmetic by
\texttt{code/verify\_r3prime.py} (receipt
\texttt{receipts/verify\_r3prime.txt}).
\end{proof}

\begin{remark}
The corrected slot count (two slots on the overlap, not three) is used
throughout; see the erratum footnote in Appendix~F. Independent
model-based referee trails (two hostile audits, two re-review rounds,
split verdict adjudicated line by line) are preserved in
\texttt{evidence/REFEREE\_R3PRIME*.md}. This is model-based review,
\emph{not} human peer review.
\end{remark}


\end{document}
